\documentclass[aspectratio=169,11pt]{beamer}

\usepackage[ngerman]{babel}
\usepackage[T1]{fontenc}
\usepackage[utf8]{inputenc}
\usepackage{lmodern}
\usepackage{amsmath,amssymb,mathtools}
\usepackage{booktabs}
\usepackage{siunitx}
\usepackage{listings}
\usepackage{pgfplots}
\pgfplotsset{compat=1.18}
\usepackage{tikz}
\usetikzlibrary{arrows.meta,positioning}

\usetheme{metropolis}

\definecolor{accent}{RGB}{0,102,153}
\metroset{titleformat=regular, numbering=fraction, progressbar=frametitle}

\title{Automatisierte Dokumentation}
\subtitle{Von der Quelle zum geprüften PDF}
\author{Erika Beispiel}
\institute{Beispiel-Institut für Technische Kommunikation}
\date{\today}

\begin{document}

\maketitle

\begin{frame}{Agenda}
  \tableofcontents
\end{frame}

\section{Motivation}

\begin{frame}{Warum Automatisierung?}
  \begin{columns}[T]
    \begin{column}{0.48\textwidth}
      \textbf{Manuell}
      \begin{itemize}
        \item Word-Dateien mit Formatverlust
        \item Verweise brechen beim Umbauen
        \item Keine reproduzierbaren Stände
      \end{itemize}
    \end{column}
    \begin{column}{0.48\textwidth}
      \textbf{Automatisiert}
      \begin{itemize}
        \item Quelltext als einzige Wahrheit
        \item Verweise und Zitate automatisch
        \item Jeder Stand per ZIP exportierbar
      \end{itemize}
    \end{column}
  \end{columns}

  \vspace{1em}
  \begin{alertblock}{Kernsatz}
    Dokumentation ist Software: Sie gehört ins Versions­system und in die
    Pipeline~\cite{knuth1984}.
  \end{alertblock}
\end{frame}

\begin{frame}{Der Build-Kreislauf}
  \centering
  \begin{tikzpicture}[
      node distance=15mm,
      st/.style={draw, rounded corners=3pt, minimum height=9mm, minimum width=25mm,
                 align=center, fill=accent!8},
      ar/.style={-{Stealth[length=2.6mm]}, thick, accent},
    ]
    \node[st] (a) {Quelltext};
    \node[st, right=of a] (b) {latexmk};
    \node[st, right=of b] (c) {PDF};
    \node[st, below=of b, fill=orange!15] (d) {Log-Auswertung};
    \draw[ar] (a) -- (b);
    \draw[ar] (b) -- (c);
    \draw[ar] (b) -- (d);
    \draw[ar, red!70!black, dashed] (d.west) -- node[above=1.5mm, font=\small]{Fix} (a.south);
  \end{tikzpicture}

  \vspace{1em}
  \footnotesize SyncTeX verbindet jede PDF-Position mit der Quellzeile.
\end{frame}

\section{Verfahren}

\begin{frame}{Metriken}
  Die Gesamtqualität ergibt sich als gewichteter Mittelwert
  \begin{equation*}
    Q \;=\; \frac{\sum_{i=1}^{n} w_i s_i}{\sum_{i=1}^{n} w_i},
    \qquad s_i \in [0,1].
  \end{equation*}

  \pause
  \begin{block}{Eigenschaften}
    \begin{itemize}
      \item Skalierungsinvariant gegenüber der Zahl der Metriken
      \item $Q = 1$ genau dann, wenn alle $s_i = 1$
      \item Empfindlich gegenüber der Gewichtung $w_i$
    \end{itemize}
  \end{block}
\end{frame}

\begin{frame}{Messergebnisse}
  \centering
  \begin{tikzpicture}
    \begin{axis}[
        width=0.86\textwidth,
        height=5.4cm,
        xlabel={Epoche},
        ylabel={$Q$},
        xmin=1, xmax=10, ymin=0.4, ymax=1.0,
        grid=major, grid style={black!15},
        legend style={at={(0.02,0.98)}, anchor=north west, draw=black!30, font=\small},
      ]
      \addplot[accent, very thick, mark=*, mark size=1.8] coordinates {
        (1,0.52) (2,0.61) (3,0.68) (4,0.73) (5,0.77)
        (6,0.80) (7,0.83) (8,0.85) (9,0.86) (10,0.871)
      };
      \addplot[red!70!black, thick, dashed] coordinates {
        (1,0.50) (2,0.56) (3,0.62) (4,0.66) (5,0.70)
        (6,0.72) (7,0.74) (8,0.75) (9,0.76) (10,0.768)
      };
      \legend{Hybrid, Regelbasiert}
    \end{axis}
  \end{tikzpicture}
\end{frame}

\section{Werkzeug}

\begin{frame}{LaTeX Studio auf einen Blick}
  \begin{description}
    \item[Live-Vorschau] PDF sofort nach dem Build, Zoom und Seitennavigation
    \item[SyncTeX] Klick ins PDF springt in den Quelltext
    \item[Autocomplete] Befehle, Umgebungen, Labels und Zitate
    \item[Problemliste] Fehler und Warnungen mit Sprungziel
    \item[Projektverwaltung] Umbenennen, Duplizieren, Export, ZIP-Import
  \end{description}

  \vspace{0.8em}
  \centering\small Tastenkürzel: \texttt{Strg+B} bauen \quad
  \texttt{Strg+P} Datei öffnen \quad \texttt{Strg+Umschalt+P} Befehlspalette
\end{frame}

\begin{frame}[fragile]{Code in Folien}
  \begin{columns}[T]
    \begin{column}{0.52\textwidth}
      \begin{lstlisting}[basicstyle=\ttfamily\tiny, frame=single,
                          language=bash, numbers=none]
latexmk -pdf -interaction=nonstopmode main.tex
# oder direkt im Editor: Strg+B
      \end{lstlisting}
    \end{column}
    \begin{column}{0.44\textwidth}
      \begin{alertblock}{Achtung}
        \texttt{shell-escape} nur für vertrauenswürdige Projekte
        aktivieren.
      \end{alertblock}
    \end{column}
  \end{columns}
\end{frame}

\begin{frame}{Zusammenfassung}
  \begin{itemize}
    \item Ein Quelltext, viele Ausgaben -- Reproduzierbarkeit inklusive
    \item Automatisierte Fehleranalyse spart Korrekturschleifen
    \item SyncTeX schließt die Lücke zwischen PDF und Quelltext
  \end{itemize}

  \vspace{1.2em}
  \centering
  \Huge\color{accent} Vielen Dank für die Aufmerksamkeit!
\end{frame}

\begin{frame}[allowframebreaks]{Literatur}
  \begingroup
  \small
  \bibliographystyle{plain}
  \bibliography{references}
  \endgroup
\end{frame}

\end{document}
